\documentclass[10pt,a4paper]{amsart}
\usepackage[margin=19mm]{geometry}
\usepackage{amsmath,amssymb,mathtools}
\usepackage[scr=rsfs,cal=euler]{mathalfa}
\usepackage{xfrac}
\usepackage[unicode,hidelinks,bookmarks=false]{hyperref}
\newcommand{\ie}{i.e.\ }
\newcommand{\cf}{cf.\ }
\newcommand{\resp}{respectively\ }
\newcommand{\Subsection}{\subsection}
\providecommand{\nobreakdash}{\nobreak\hbox{-}}
\setlength{\emergencystretch}{2em}
\multlinegap0pt
\newcommand{\varw}{z}
\newcommand{\varz}{w}
\usepackage{amssymb}
\newtheorem{lemma}{Lemma}
\newcommand{\IC}{\mathbb{C}}
\newcommand{\IN}{\mathbb{N}}
\newcommand{\pco}{\mathrm{PCO}}
\renewcommand{\subset}{\subseteq}
\title{Dynamical Canonical Heights and Finite Trees}
\author{Philipp Habegger, Harry Schmidt}
\date{Source: arXiv:2608.00774v1 (CC BY 4.0)}
\begin{document}
\maketitle

\noindent\emph{Source and licence.} Dynamical Canonical Heights and Finite Trees. Version arXiv:2608.00774v1 (CC BY 4.0), distributed by arXiv under the Creative Commons Attribution 4.0 International licence, http://creativecommons.org/licenses/by/4.0/, as recorded in the arXiv licence metadata for this exact version and shown on the abstract page https://arxiv.org/abs/2608.00774v1. Full licence notice: \texttt{licenses/arxiv-2608.00774-CC-BY-4.0.txt}.\par\smallskip

\noindent\textit{Editorial scope.} Counted unit: the article's lemma lem:fninvuniform\_new1 (source0.tex lines 6571--6594). The article's complete original proof of it is delivered with it (source0.tex lines 6595--6798, one \texttt{proof} environment of 204 lines). No mathematics is written, repaired or re-derived: the statement and the proof are the article's own lines, in the article's own notation, and no retained line is substituted or re-flowed.  Retained uncounted as the source-local context the proof needs: lines 6226--6236; lines 6371--6378; lines 6463--6476; lines 6491--6499; lines 6528--6536.  Nothing was omitted from any retained span, so the package carries no omission record.  No retained line contains a citation, so the wrapper carries no bibliography and no bibliography entry.  No retained line contains a figure environment, an image-inclusion command or drawing code, so no image is delivered and the package declares no asset dependency.  Editorial formatting only, in addition: an AMS \texttt{amsart} preamble that restates the article's own declarations and macros used by the retained lines, namely the environments \texttt{lemma} on separate counters, together with the standard packages those lines need.  The retained environments print in this document as Lemma 1 in order of appearance, whereas the article prints the same items with the numbering of its own full document.  The complete native source of the article as distributed in the arXiv e-print for this exact version is bundled unchanged as \texttt{sources/206566/source0.tex}.
\begin{lemma}
  \label{lem:hyperboliccharacterization}
  Let $f\in\IC[T]$ be of degree at least $2$.  Then the following are
  equivalent.
  \begin{enumerate}
  \item [(i)] The polynomial $f$ is hyperbolic.
  \item[(ii)] The closure of $\pco_{\ge 0}(f)$ in $\IC$ is disjoint from $J_f$.
  \item[(iii)] There exists an integer $m\ge 0$ such that
    closure of $\pco_{\ge m}(f)$ in $\IC$ is disjoint from $J_f$.
  \end{enumerate}
\end{lemma}

\begin{lemma}
  \label{lem:hyperbolic2}
  Suppose $f$ is as above and hyperbolic.  
  There exist $\epsilon>0,r\ge 1,$ and $\lambda>1$ that depend
  only on $f$ with the following property. For all $z\in\IC$ and all
  $s\in\IN_0$ we have $\mathrm{dist}(f^{(s)}(z),J_f) \ge \epsilon
  \min\{1,\lambda^s\, \mathrm{dist}(z,J_f)\}^r$.
\end{lemma}

\begin{lemma}
  \label{lem:hyperboliccovering}  
  Suppose $f\in\IC[T]$ has degree at least $2$ and is hyperbolic.
  For all sufficiently small $\epsilon > 0$ the following holds.
  Let $w\in J_f$ and $s\in\IN_0$.
  \begin{enumerate}
  \item [(i)]
    The restriction of $f^{(s)}$ to $(f^{(s)})^{-1}(B_\epsilon(w))$ is a
    covering map onto $B_\epsilon(w)$. 
  \item[(ii)] Let $V$ be a connected component of
    $(f^{(s)})^{-1}(B_\epsilon(w))$. Then $f^{(s)}|_V\colon
    V\rightarrow B_\epsilon(w)$ is a biholomorphism.
  \end{enumerate}
\end{lemma}

\begin{lemma}
  \label{lem:koebequarterapp}
  Let $U\subset\IC$ be a domain and let $g\colon U\rightarrow\IC$ be
  an injective and holomorphic map. Suppose $w,w'\in U$  such that
  $\overline{B_{|w-w'|}(w)}\subset U$. Then
  \begin{equation*}
    |g(w)-g(w')|\ge \frac{|g'(w)|}{4}|w-w'|.
  \end{equation*}
\end{lemma}

\begin{lemma}
  \label{lem:elementary}
  Let $z_1,\ldots,z_s\in\IC$ and $\delta_1,\ldots,\delta_s \in (0,1/2)$
  with $||z_i|-1|\le \delta_i$. Set
  $\delta=\sum_{i=1}^s\delta_i$, then 
  \begin{equation*}
    \bigl||z_1\cdots z_s|-1\bigr| \le e^{2\delta}-1.
  \end{equation*}  
\end{lemma}

\begin{lemma}
  \label{lem:fninvuniform_new1}
  Suppose $f\in\IC[T]$ is of degree at least $2$ and
  hyperbolic. The following holds for all $\epsilon > 0$  sufficiently small in terms of
  $f$. Let
  $s\in\IN_0$, $\widehat w\in J_{f}$, and let
  $V$ be a connected component of $(f^{(s)})^{-1}(B_{\epsilon}(\widehat{w}))$.
  Then $f^{(s)}|_V\colon V\rightarrow B_{\epsilon}(\widehat{w})$ is
  a biholomorphism.
  Let $g\colon
  B_{\epsilon}(\widehat{\varz})\rightarrow V$ denote the inverse map and
  $\widehat z = g(\widehat w)$. The following hold true.
  \begin{enumerate}
  \item[(i)] 
    We have $\frac 12 |g'(\widehat w)|\le |g'(w)|\le 2|g'(\widehat w)|$ for all $w\in
    B_{\epsilon}(\widehat w)$. 
  \item[(ii)]  
    We have ${f^{(s)}}'(\widehat{\varw})\not=0$ and
    \begin{equation*}
      \frac 12 |z-z'| \le \frac{|{f^{(s)}}(z)-{f^{(s)}}(z')|}{|{f^{(s)}}'(\widehat{\varw})|} \le 8 |z-z'|
    \end{equation*}
    for all $z,z'\in V$. 
  \end{enumerate}
\end{lemma}

\begin{proof}
  If $s=0$, then $f^{(s)}$ is the identity and so is $g$. In this
  case, the lemma is trivial. So we assume $s\ge 1$. 
  We let $\epsilon \le 1$ be small enough so that 
  Lemma~\ref{lem:hyperboliccovering} holds.
  The first claim, on $f^{(s)}|_V$, holds by the said lemma.
  
  In the proof let
  $g\colon B_{\epsilon}(\widehat w)\rightarrow V$ be the inverse of
  $f^{(s)}|_V\colon V\rightarrow B_{\epsilon}(\widehat w)$.

  For brevity, we write $U = B_{\epsilon}(\widehat w)$. Let us define
  $g_k = f^{(s-k)}\circ g\colon B_{\epsilon}(\widehat w)\rightarrow \IC$
  for all $k\in \{1,\ldots,s\}$. In particular, $g_s=g$. It is
  convenient to define $g_0$ to be the identity.

  We will shrink $\epsilon\in (0,1]$ several times. The values of
  $\lambda,c,l,c_1,c_2,\ldots$ are positive and fixed in the proof. They
  will depend only on $f$ but not on $\widehat w,s,V$ or the branch
  $g$ of $(f^{(s)})^{-1}$.


  We first prove that there exists $\lambda > 1$
  that depends only on $f$
  such that
  \begin{equation}
    \label{eq:fninvuniform}
    \mathrm{dist}(g(\varz),J_f) \le c \lambda^{-s}
    \quad\text{and}\quad
    |g'(\varz)|\le c\lambda^{-s}
    \quad\text{for all}\quad w\in U. 
  \end{equation}


  Say $\varz\in U$ and $k\in \{0,\ldots,s\}$. By Lemma~\ref{lem:hyperbolic2}
  applied to $g_k(\varz)$ and $k$ we find
  \begin{equation*}
    \mathrm{dist}(\varz,J_f) \ge c_1^{-1}
    \min\{1,\lambda^k \mathrm{dist}(g_k(\varz),J_f)\}^r;   
  \end{equation*}
  here $c_1>0, \lambda > 1,$
  and $r\ge 1$ depend only on $f$.
  We may assume $\epsilon \le 1/c_1$. 
  As $\mathrm{dist}(\varz,J_f)<\epsilon\le 1/c_1$ we get
  $\epsilon >  c_1^{-1}
  \lambda^{rk} \mathrm{dist}(g_k(\varz),J_f)^r$.
  So 
  \begin{equation}
    \label{eq:gkapproaching}
    \mathrm{dist}(g_k(\varz),J_f) \le (c_1\epsilon)^{1/r} \lambda^{-k} < (c_1\epsilon)^{1/r}
    \quad\text{for all } \varz\in U\text{ and all }k\in
    \{0,\ldots,s\}. 
  \end{equation}
  Recall that $g_s=g$.
  The first inequality in 
  (\ref{eq:fninvuniform}) follows with
  $c = (c_1\epsilon)^{1/r}$.

  As $f$ is hyperbolic, there exist $l\in\IN$ and $\mu > 1$ with
  $|F'(\varw)|\ge \mu$ for all $\varw$ sufficiently close to $J_f$ and where
  $F = f^{(l)}$. 
  This holds in particular for all  $z\in g_k(U)$ and all $k\in
  \{0,\ldots,s\}$ by
  (\ref{eq:gkapproaching}) as we may assume $\epsilon$ is small in
  terms of $f$.

  We continue to assume $k\in\{0,\ldots,s\}$ and 
  write $k=ml+t$ for some $m\in\IN_0$
  and $t\in \{0,\ldots,l-1\}$. Thus
  $F^{(m)}\circ f^{(t)}\circ g_k$  is the identity on $U$.
  Let $\varz\in U$ be arbitrary, the chain rule yields
  \begin{equation}
    \label{eq:chainrulehyperbolic}
    1 =    {f^{(t)}}'(g_k(\varz)) g'_k(\varz)\prod_{j=1}^m
      F'\bigl(F^{(m-j)}(f^{(t)}(g_k(\varz)))\bigr).
  \end{equation}
  If $j\in \{1,\ldots,m\}$, then
  $F^{(m-j)}(f^{(t)}(g_k(\varz))) = f^{(k-jl)}(g_{k}(\varz)) =
  g_{jl}(\varz)$. So 
  \begin{alignat}1
    \label{eq:chainrulehyperbolic2}
    1 &=
    \left|      {f^{(t)}}'(g_k(\varz)) g'_k(\varz)\prod_{j=1}^m F'\bigl(g_{jl}(w)\bigr)
\right|
    \ge    |      {f^{(t)}}'(g_k(\varz))| |g_k'(\varz)|\mu^m.
  \end{alignat}

  The derivative ${f^{(t)}}'$ does not vanish on $J_f$ for any $t\in
  \{0,\ldots,l\}$.
  Indeed, if
  ${f^{(t)}}'(w_0)=0$ then $f^{(t)}(w_0) \in \pco_{\ge 1}(f)$ by the
  chain rule.
  So $w_0\not\in J_f$ by Lemma~\ref{lem:hyperboliccharacterization}
  and since $f$ is hyperbolic. 
  So $|{f^{(t)}}'|$ is bounded from below
  by a positive value on some neighborhood of the compact set $J_f$. For
  $\epsilon>0$ small enough in terms of $f$, (\ref{eq:gkapproaching}) guarantees
  that $g_k(w)$ is sufficiently close to $J_f$. Thus
  \begin{equation*}
    |{f^{(t)}}'(g_k(w))|\ge c_2
    \text{ for all } t\in
    \{0,\ldots,l\},\text{ all }k\in\{0,\ldots,s\},\text{ and all }w\in U
  \end{equation*}
  where $c_2>0$ depends only on $f$ and $l$. But $l$ was fixed in
  terms of $f$, so $c_2$ depends only on $f$. 
  In particular, and this will be useful later on, $f'(g_k(w))\not=0$ for all $w\in U$.
  From (\ref{eq:chainrulehyperbolic2}) 
  we deduce
  \begin{equation*}
    |g_k'(w)|\le c_3 \mu^{-m}\quad\text{for all}\quad w\in U
    \text{ and all } k\in \{0,\ldots,s\}
  \end{equation*}
  where $c_3>0$ depends only on $f$. 
  As $k< ml+l$ we conclude $m> k/l-1$ and
  \begin{equation}
    \label{eq:gkprime}
    |g_k'(\varz)|\le c_4  \mu^{-k/l}\text{ for all } \varz\in U
    \text{ and all } k\in \{0,\ldots,s\}.
  \end{equation}
  Again $c_4>0$ and also $\mu>1$ and $l$ 
  are independent of $s,k,$ and the
  choice of branch $g$ of $(f^{(s)})^{-1}$.
  
  So the second bound in the claim (\ref{eq:fninvuniform}) holds after
  adjusting $\lambda$ in the case $k=s$.

  The claim is now proved. We now turn to conclusions (i) and (ii). 

  Let $w,w'$ lie in the disk $U$. 
  Then $|g_k(w)-g_k(w')| \le \sup_{w''\in L} |g_k'(w'')|
  |w- w'|$ where $L\subset U$ is the line segment connecting
  $w$ to $w'$; just integrate $g'_k$ along $L$. 

  So $|g_k(\varz)-g_k(\widehat{\varz})|\le 2c_4
  \mu^{-k/l} |\varz- \widehat{\varz}|$ for all $\varz\in U$
  using (\ref{eq:gkprime}). 
  By (\ref{eq:gkapproaching}) and compactness this implies
  $|f'(g_k(\varz))-f'(g_k( \widehat{\varz}))| \le c_5 \mu^{-k/l}|\varz- \widehat{\varz}|$.
  We saw above that $|f'(g_k(w))|\ge c_2$. Taking $c_6=c_5/c_2$ yields
  \begin{equation}
    \label{eq:1minusfprimef}
    \left|1-\frac{f'(g_k( \widehat{\varz}))}{f'(g_k(\varz))}\right|\le c_6
    \mu^{-k/l}|\varz- \widehat{\varz}|\quad\text{for all}\quad k\in
    \{0,\ldots,s\}\quad\text{and all}\quad w\in U.
  \end{equation}
  The chain rule implies $1 = g'(w) {f^{(s)}}'(g(w)) = g'(w)
  f'(f^{(s-1)}(g(w)))\cdots f'(g(w))$. We use the
  definition of $g_k$ and find
  $1 = g'(w)f'(g_1(w)) \cdots f'(g_s(w))$ on $U$.
  This holds in particular for $\widehat w$.
  We take the
  quotient at $\varz$ and $ \widehat{\varz}$ and find  
  \begin{equation*}
    \frac{g'(\varz)}{g'( \widehat{\varz})} = \frac{f'(g_1( \widehat{\varz}))}{f'(g_1(\varz))} \cdots \frac{f'(g_s( \widehat{\varz}))}{f'(g_s(\varz))}.
  \end{equation*}

  All factors on the right are near $1$ by   (\ref{eq:1minusfprimef})
  and for $\epsilon$ small.
  Indeed, as $|w-\widehat w |<\epsilon$ we may assume
  $c_6\mu^{-k/l}|w-\widehat w|< 1/2$.
  Thus $\sum_{k=1}^{s} c_6\mu^{-k/l}|w-\widehat w|
  \le {c_6}(1-\mu^{-1/l})^{-1}|w-\widehat w|\le c_7 |w-\widehat w|$. We
  refer to Lemma~\ref{lem:elementary} and obtain
  \begin{equation*}
    \left|\left|\frac{g'(\varz)}{g'( \widehat{\varz})}\right|-1\right|
    \le e^{2c_7|w-\widehat w|}-1\le c_8|\varz- \widehat{\varz}|.
  \end{equation*}
  We may assume $c_8\epsilon \le 1/2$.
  So 
  $$\frac 12 \le 1-c_8|\varz-\widehat{\varz}|\le
  \left|\frac{g'(\varz)}{g'(\widehat{\varz})}\right| \le
  1+c_8|\varz-\widehat{\varz}|\le 2$$
  for all $\varz\in U$.
  So $|g'(\varz)|\le 2  |g'(\widehat{\varz})|$ and
  $|g'(\widehat\varz)|\le 2 |g'({\varz})|$.
  This implies part (i).

    Let $z,z'\in V$ be as in (ii).
  We set $w=f^{(s)}(z)$ and $w'=f^{(s)}(z')$, both elements of $U$. 
  Recall that 
  $|g(w)-g(w')|\le  \sup_{w''\in L} |g'(w'')| |w-w'|$.
  So 
  $|g(w)-g(w')|\le 2|g'(\widehat{w})||w-w'|$ by part (i). 
  Now $g'(\widehat{w}){f^{(s)}}'(\widehat z)=1$ by the chain rule. So
  $|g(w)-g(w')|
  \le 2|{f^{(s)}}'(\widehat{z})|^{-1}|w-w'|$. 
  This completes the first inequality in (ii).

  For the second inequality we assume $z,z' \in
  g(B_{\epsilon/3}(\widehat w))$ and set
  $w=f^{(s)}(z)$ and $w'=f^{(s)}(z')$. So $|w-w'|<2\epsilon/3$ and the
  closed disk $\overline{B_{2\epsilon/3}(w)}$ lies entirely inside
  $B_\epsilon(\widehat w)$.
  By Lemma~\ref{lem:koebequarterapp} and part (i) we find
  $|g(w)-g(w')|\ge |g'(w)||w-w'|/4 \ge |g'(\widehat w)| |w-w'|/8$.
  As we have seen above, $g'(\widehat w) =1/{f^{(s)}}'(\widehat
  z)$.  Moreover, $g(w)=z$ and $g(w')=z'$. So
  we conclude the second bound in (ii) for all pairs in
  $g(B_{\epsilon/3}(\widehat w))$. We shrink $\epsilon$ a
  final time to get the second inequality of part (ii)
  on $g(B_\epsilon(\widehat w))$. The
  already proven
  estimates in (i) and (ii) remain true.
\end{proof}

\end{document}
